\section{Archon: búsqueda de Inference-Time Architecture}

\subsection{La arquitectura de inferencia vista como algo que se puede buscar}

Archon ya no elige a mano una sola técnica, sino que trata las inference-time operations como capas de una red neuronal y busca su combinación, orden, procedencia de modelos y número de invocaciones.

\videofigure{figures/archon-framework.jpg}{Archon busca una arquitectura de inferencia optimizada según el benchmark objetivo y los modelos disponibles.}{00:45:47--00:48:36}

Su entrada incluye la tarea objetivo, los modelos candidatos, las operaciones disponibles y las restricciones de costo; su salida es un grafo de ejecución. Por ejemplo, varios generator muestrean en paralelo, un critic identifica defectos, un ranker selecciona candidatos, un fuser sintetiza la respuesta, y luego se realiza otra ronda de crítica y fusión.

\subsection{Operaciones componibles}

\videofigure{figures/archon-operations.jpg}{Definición, entradas, salidas y costo de las inference-time operation más comunes en Archon.}{00:48:36--00:51:20}

\begin{table}[H]
\centering
\small
\begin{tabular}{p{2.8cm}p{5.1cm}p{5.7cm}}
\toprule
\textbf{Operación} & \textbf{Función} & \textbf{Riesgo principal} \\
\midrule
Generator & Produce uno o más candidatos & Falta de diversidad, errores repetidos \\
Critic & Señala defectos y omisiones de los candidatos & Crítica imprecisa o centrada solo en la forma superficial \\
Ranker & Ordena los candidatos & Sesgo de posición, sesgo de longitud, sesgo del judge \\
Fuser & Sintetiza las partes complementarias de varios candidatos & Fusionar también los errores \\
Verifier / Tests & Verifica restricciones y corrección & Cobertura de pruebas insuficiente, reward hacking \\
\bottomrule
\end{tabular}
\end{table}

\subsection{Por qué la fusion puede superar a la oracle selection}

La oracle selection solo puede elegir una respuesta completa entre los candidatos existentes; el fuser puede extraer información complementaria de varios candidatos incompletos y generar una respuesta nueva que no existía en el conjunto original de candidatos.

\videofigure{figures/fusion-results.jpg}{Rank-then-fuse supera, en algunas tareas, a la simple selección de un solo candidato.}{00:51:20--00:54:58}

\begin{warningbox}{"Superar al oracle" no contradice la definición}
Aquí el oracle solo selecciona la respuesta correcta entre los candidatos originales; la fusion ejecuta nuevas invocaciones al modelo y genera nuevos candidatos, por lo que cambia el espacio de candidatos. No es un selector más potente, sino un paso de generación adicional.
\end{warningbox}

\subsection{Pruebas unitarias en lenguaje natural}

Para las tareas que carecen de un ground truth ejecutable, Archon hace que el modelo genere pruebas o restricciones en lenguaje natural, y luego las usa para verificar las respuestas candidatas. Esto equivale a descomponer primero un objetivo ambiguo en especificaciones verificables localmente.

\videofigure{figures/natural-language-tests.jpg}{Los unit tests en lenguaje natural convierten los requisitos de la tarea en condiciones verificables.}{00:55:30--00:56:24}

Esta práctica mejora la capacidad de verificación estructurada, pero, como las pruebas las genera el modelo, aún pueden omitirse requisitos clave. La forma más segura de usarla es combinar los requisitos explícitos del usuario, las reglas del dominio y las pruebas generadas por el modelo, en lugar de depender por completo de especificaciones autogeneradas.

\subsection{Arquitecturas de razonamiento profundo y búsqueda de arquitectura}

\videofigure{figures/archon-architectures.jpg}{Estructura multicapa generator--critic--ranker--fuser encontrada por Archon.}{00:56:24--00:59:32}

El espacio de arquitecturas crece de forma explosiva con el número de capas, el tipo de operación, la elección de modelos y las combinaciones de hiperparámetros. Archon usa optimización bayesiana para predecir, con menos ensayos, qué configuraciones tienen más probabilidad de mejorar la métrica objetivo.

\videofigure{figures/bayesian-search.jpg}{La optimización bayesiana descubre arquitecturas de alto rendimiento con más eficiencia que la búsqueda aleatoria o voraz.}{00:59:32--01:01:20}

Si la configuración es $a$, la función objetivo puede escribirse como una utilidad multiobjetivo:

$$
J(a)=\mathrm{Quality}(a)-\lambda_c\mathrm{Cost}(a)-\lambda_l\mathrm{Latency}(a).
$$

\begin{itemize}
    \item $\mathrm{Quality}(a)$: precisión o puntaje de preferencia en el benchmark objetivo;
    \item $\mathrm{Cost}(a)$: costo de invocaciones al modelo y de tokens;
    \item $\mathrm{Latency}(a)$: latencia de la ruta crítica;
    \item $\lambda_c,\lambda_l$: pesos que el escenario de despliegue asigna al costo y a la latencia.
\end{itemize}

\videofigure{figures/archon-results.jpg}{Archon logra mejoras consistentes en pools de modelos de código abierto, cerrados y mixtos.}{01:01:20--01:03:02}

\subsection{Resumen de la sección}

Archon eleva el "prompt engineering" a una búsqueda de arquitectura de sistema: no solo optimiza un prompt, sino el flujo de datos, el flujo de control y el costo entre varias invocaciones al modelo. Su limitación es que el resultado de la búsqueda puede sobreajustarse al benchmark, y que la latencia, la interpretabilidad y el costo de mantenimiento de las arquitecturas complejas aumentan de forma considerable.
